\documentclass[a4paper,11pt]{article}

\usepackage[utf8]{inputenc}
\usepackage[T1]{fontenc}
\usepackage[german]{babel}
\usepackage[a4paper,margin=2.5cm]{geometry}
\usepackage{enumitem}
\usepackage{titlesec}
\usepackage{longtable}
\usepackage{setspace}
\usepackage{parskip}

\begin{document}

% ---------------------------
% Deckblatt
% ---------------------------

\begin{titlepage}
    \centering
    
    \vspace*{3cm}
    
    {\Huge \textbf{Meetingprotokoll}}\\[0.5cm]
    {\Large Statusbericht zur Projektarbeit}\\[2cm]
    
    \rule{12cm}{0.5pt}\\[0.5cm]
    
    {\Large Projektarbeit}\\[0.3cm]
    {\large 12. Mai 2026}
    
    \rule{12cm}{0.5pt}
    
    \vfill
    
    {\large THWS / MINOVA}\\
    
\end{titlepage}

% ---------------------------
% Inhaltsverzeichnis
% ---------------------------

\tableofcontents
\newpage

% ---------------------------
% Inhalt
% ---------------------------

\section{Allgemeine Informationen}

\begin{tabular}{ll}
\textbf{Datum:} & 12.05.2026 \\
\textbf{Art des Meetings:} & Projektmeeting / Statusbesprechung \\
\textbf{Nächstes Meeting:} & 18.05.2026 um 17:30 Uhr \\
\end{tabular}

\vspace{1cm}

\section{Projektstatus}

Der aktuelle Stand des Projekts wurde im Meeting besprochen und dokumentiert. 
Die grundlegende Backend-Struktur wurde bereits implementiert. 
Außerdem wurden erste konzeptionelle und organisatorische Bestandteile des Projekts vorbereitet.

\subsection{Aktuell abgeschlossene bzw. bearbeitete Aufgaben}

\begin{itemize}[leftmargin=1.2cm]

    \item \textbf{Backend}
    \begin{itemize}
        \item Erster Datenbankentwurf erstellt
        \item Erste grundlegende BPMN-Diagramme erstellt (in Zusammenarbeit mit dem Backend)
        \item Strukturdiagramm aus Sicht des Backends vorbereitet
        \item Backend-Grundstruktur implementiert
    \end{itemize}

    \item \textbf{Dokumentation}
    \begin{itemize}
        \item Erstes E-Mail-Beispiel erstellt
        \item Lastenheft und Pflichtenheft ausgearbeitet
        \item Projektplan erstellt
    \end{itemize}

\end{itemize}

\vspace{0.5cm}

\subsection{Verschobene Aufgaben}

Die Arbeiten am Frontend wurden aufgrund des ITW-Moduls vorübergehend verschoben.

\vspace{0.5cm}

\section{Besprochene Punkte}

Während des Meetings wurden folgende organisatorische und technische Themen besprochen:

\begin{itemize}[leftmargin=1.2cm]
    \item Für jeden Projektermin muss ein offizielles Meetingprotokoll erstellt werden.
    \item Der Präsentationstermin steht aktuell noch nicht fest.
    \item Die Stundenaufstellung wird kontinuierlich dokumentiert.
\end{itemize}

\vspace{0.5cm}

\section{Nächste Schritte (nächste 2 Wochen)}

\begin{itemize}[leftmargin=1.2cm]
    \item Entwicklung eines funktionierenden Prototyps
    \item Beginn der Frontend-Grundstruktur (Frontend-Skelett)
    \item Erweiterung der bestehenden Backend-Funktionalitäten
    \item Besprechung mit dem Auftraggeber bezüglich der weiteren Projektplanung
\end{itemize}


\vspace{0.5cm}

\section{Nächste Projektphase}

Die nächste Projektphase konzentriert sich auf:

\begin{itemize}[leftmargin=1.2cm]
    \item Definition der Systemarchitektur
    \item Planung und Beschreibung der Schnittstellen
    \item Technische Abstimmung zwischen Frontend und Backend
\end{itemize}

\vspace{0.5cm}

\section{Nächste Meilensteine}

\begin{itemize}[leftmargin=1.2cm]
    \item Systemarchitektur definieren
    \item DISPO-Schnittstelle definieren
    \item Backend-Grundstruktur implementieren (bereits abgeschlossen)
\end{itemize}

\vspace{1cm}

\section{Zusammenfassung}

Das Projekt befindet sich aktuell in einer stabilen frühen Entwicklungsphase. 
Die grundlegenden technischen und organisatorischen Strukturen wurden vorbereitet beziehungsweise bereits implementiert. 
In den kommenden Wochen liegt der Fokus insbesondere auf der Weiterentwicklung der Endpunkte, dem Start der Frontend-Entwicklung sowie der Definition der Systemarchitektur und der Schnittstellen.

\end{document}